\section{Methods}
\label{sec:methods}

\subsection*{Agent skeleton and task domain}

All experiments use a CodeAct-style coding agent executing ScienceAgentBench (SAB) tasks from the verified snapshot \texttt{sab-verified-c26e151} \cite{chen2024scienceagentbench}. Each run receives one task instance whose instruction requires multi-step scientific data analysis in a Python sandbox; the agent interleaves reasoning and code execution for up to 12 steps, with a 300-second per-step execution timeout and an 8,192-token generation cap, at temperature 0 throughout. Task success is adjudicated exclusively by the benchmark's official evaluation script for the instance --- a programmatic checker over files the agent produces --- so the outcome label Y $\in$ {0,1} contains no model-judge component. Every run executes in a freshly created sandbox working directory; directory permissions restrict the agent to its own workspace, and the experiment harness runs under an unprivileged account. The frontier policy is GPT-5.5 (snapshot \texttt{gpt-5.5-2026-04-23}) accessed through an OpenAI-compatible relay; the frontier model identity, snapshot and fallback snapshot are recorded in the run manifest of every run, and the full trajectory (prompts, code, observations, step timings) is persisted for each run.

\subsection*{Step-level counterfactual replacement protocol}

A \emph{parent trajectory} is a complete successful frontier run on a task instance. The protocol freezes the parent --- prompts, produced artifacts, sandbox state and the step index at which each artifact appeared --- and re-executes the trajectory from a chosen step t under one of four \emph{replacement arms}: (i) \texttt{frontier-resample}, where the frontier model itself regenerates step t, controlling for the stochasticity of re-execution; and three \emph{edit-model} arms in which step t is regenerated by a cheaper policy --- GPT-5.4-mini (same model family as the frontier), DeepSeek-V3.2 (cross-family), and Qwen3-8B served locally from bf16 weights via vLLM (a 14B AWQ quantised variant is used in the scale-replication experiment). In every arm the replacement receives exactly the prompt and observation state the parent's step t received, and after step t the frontier policy resumes execution for the remainder of the trajectory, so that post-replacement competence is held constant and any success-rate difference is attributable to the replaced step. Each (point, arm) cell is replicated (minimum two runs; every one of the 38 analysed points has at least two valid runs in all four arms after targeted backfilling), and a run counts as valid only if it reaches the task's evaluation stage inside the step budget; budget-exhaustion is separately audited by a protocol-retry experiment (below).

\subsection*{Signal and control points}

Thirty-eight step instances (points) were analysed: 6 \emph{signal points}, identified in a first-pass comparison as step instances where replacement arms separated visibly from the frontier resample baseline, and 32 \emph{control points} sampled across the remaining step inventory without reference to replacement outcomes. Because the signal points were selected on the same data that later analyses summarise, all headline effect estimates use difference-in-differences against matched controls rather than raw signal-point rates, so that the reported quantity is the step-specific component of damage beyond what the same arms show on ordinary steps; the control set also anchors the prediction analysis, where points enter uniformly regardless of origin. Point identity, task instance, step index and parent-run identifier are recorded in a versioned inventory table; two late-labelled points (initially missing runs in one or more arms) were backfilled with additional controlled runs before unblinding, and the inventory version used for all analyses contains 772 runs in total (767 experiment runs plus targeted infrastructure runs reported separately in the reproducibility table).

\subsection*{Full-step scan and mechanism interventions}

To obtain a dense profile rather than a sampled one, two representative parent trajectories were swept across \emph{every} step and all four arms under the identical protocol --- task t40 (five steps before its first deliverable-producing step) and task t60 (nine steps) --- yielding 193 sandbox-valid runs. Three interventions were then applied at the signal points. \emph{Structured feedback}: when a replacement run fails at a signal step, the edit model is re-prompted with a structured diagnosis (error class, expected output contract, failing assertion) and re-executes the step; the bare-retry control re-prompts without diagnosis. \emph{Suffix takeover}: the cheaper model's edited artifact is spliced into the frontier trajectory and the frontier model continues from the splice point --- testing whether step-level dependence can be repaired by local substitution plus frontier recovery. \emph{Scale replication}: the open-weights arm is re-run at Qwen3-14B (AWQ) on all six signal points to test whether step-level failures dissolve with model scale. Finally, a \emph{protocol retry} control re-executes runs that exhausted the step budget under an extended observation window, to verify that budget exhaustion reflects the observation window rather than a property of replaced steps.

\subsection*{Blind annotation of verifiability features}

For each of the 38 points, a materials-generation script produced a \emph{reading pack} containing the task instruction verbatim, the full step source code (truncation ceiling 25,000 characters; no point exceeded it), and the first 1,200 characters of the step's observation. The pack contains no message-history summary, no success labels, no model or arm information, and no per-point outcome of any kind; the script embeds machine checks (verified file content rather than file names, mandatory field presence) that were added after a generation defect was caught and corrected. Two annotators independently scored seven features per point from a frozen codebook (v1.3): F1 objective verifiability (does the step's output have an objective correctness standard), F2 validator availability (operationalised: does the step write to the benchmark's deliverable directory, produce a locally checkable intermediate artifact, or produce nothing), F3 feedback availability (was error information available at execution time), F4 feedback informativeness (how informative was the available feedback), F5 state dependence (dependence on prior-step conventions visible in code and observations), F6 recoverability (a priori structural recoverability from the step's own information), and F7 output type (categorical). Annotator A is an AI assistant operating under a declared contamination boundary (it contributed to materials drafting and holds step-type priors but never saw per-point replacement outcomes; its residual-risk disclosure is versioned in the repository); annotator B is the project lead, who received only the reading packs, the codebook and a blank template. Agreement was scored as weighted Cohen's $\kappa$ (quadratic weights) for F1--F6 and unweighted $\kappa$ for F7; features falling below 0.6 entered a pre-registered two-round clarification-and-reannotation procedure, and remaining disagreements were adjudicated on a {0, 0.5, 1, 1.5, 2} value grid (differences of at most one level averaged; larger differences resolved by consensus). The merged 38 $\times$ 7 matrix (128 cells by agreement, 81 by one-level averaging, 17 by adjudication, 2 by annotator-A default) was committed with a timestamp before any per-point replacement outcome was examined, and the commit hash precedes the unblinding commit in the repository history.

\subsection*{Prediction protocol}

The prediction target is the per-arm success rate of each (point, arm) cell --- 134 valid observations over 38 points --- and the pre-registered hypothesis H2 states that verifiability features carry predictive signal beyond controls. Controls C comprise step index, remaining budget, cliff-stratum membership, deliverable-role membership and arm identity; features V comprise F1--F6 (F7 is dropped as a constant). Model M0 is ridge regression on C alone; M1 adds V; the regularisation strength is fixed at $\alpha$ = 1.0 before any evaluation, with no tuning on held-out data. Validation is leave-one-task-out cross-validation over the 14 task instances, guaranteeing that no two observations from the same parent trajectory ever straddle a fold boundary. Metrics are pooled RMSE against per-arm success rates, Spearman rank correlation, and AUROC of success events (Y $\geq$ 0.5). The pre-registered decision rule declares support for H2 if and only if $\Delta$RMSE $\leq$ 0 and $\Delta$$\rho$ $\geq$ 0.1 (M1 relative to M0); both were met. A gradient-boosted-tree specification is run as a sensitivity check under the identical split and metrics, and per-arm decomposition and gain attribution (which points the feature model helps or hurts most) are reported alongside.

\subsection*{Held-out validation}

After the primary analysis was frozen, the seven scan code-step instances outside the 38-point selection were designated as a held-out set (the two scan instances that are final, code-free steps were excluded so that the held-out composition matches the annotated set). Their verifiability features were coded from freshly generated, leak-checked reading materials by the same two-annotator protocol, with the role control set to the conservative neutral value for all arms and the stratum control taken from the frozen task-stratum matrix; both mappings were fixed before predictions were unblinded. Two training configurations were evaluated --- all 38 points (instance-level) and the 32 points remaining after removing every scan-task point (task-level) --- each fitted with the frozen ridge specification and evaluated on the 28 held-out observations by the same metrics, with a 20,000-fold permutation test attached to the rank correlation. Because the held-out outcomes appear in the scan reports, this validation is a held-out reanalysis rather than a prospective test, and the annotation-independence caveat for annotator A is correspondingly wider; both facts are stated wherever the result is cited.

\subsection*{True cross-domain validation}

As a fully prospective test, two Computational Chemistry tasks were drawn from the same verified benchmark pool --- a domain absent from all 14 training tasks --- and carried through the complete protocol from scratch: new frontier parents were run (two tasks passed the fidelity and evaluation gates; four additional parent attempts failed and were discarded), all three code steps of each parent were pre-registered as replacement points under a rule fixed before any counterfactual run, and the full 4-arm $\times$ 2-replication grid was executed (46 of 48 runs landed; repeated relay timeouts left two deepseek cells at one replication, so every point--arm cell retains at least one observation). The six new points were annotated from freshly generated, leak-checked materials by the same two-annotator protocol (34 of 36 feature cells in exact consensus, 2 one-level disagreements averaged, none requiring adjudication); because the new tasks carry no stability-matrix entry and no role labels, the stratum and role controls were fixed at their neutral values before unblinding, matching the conservative choice of the held-out reanalysis. Three training configurations were evaluated --- the 38 pre-registered points, the 32 points remaining after removing the scan tasks, and 45 points adding the seven previously held-out scan instances --- each fitted with the frozen ridge specification, evaluated on the 24 held-out observations, and accompanied by 20,000-fold permutation tests on the rank correlation.

\subsection*{Statistics}

Pairwise arm comparisons at individual points use Fisher's exact tests with Benjamini-Hochberg false-discovery-rate correction across the family of step-arm tests (42 in the scan snapshot); pooled signal-point contrasts use stratified Cochran-Mantel-Haenszel tests with the capability-dependent points as strata; success-rate point estimates carry Wilson 95\% confidence intervals throughout. The power situation is reported rather than obscured: with 38 points and per-cell replication of two, no individual step-arm contrast in the full-step scan survives FDR correction at q < 0.10, and the confirmatory claims rest on the aggregations specified above; the prediction result is directional against the pre-registered target of $\geq$ 60 annotated points. All tests are two-sided, exact p-values are reported, and no post-hoc exclusions were performed; the one protocol deviation in the scan (a snapshot boundary described in the reproducibility table) is documented with both calibres side by side.

\subsection*{Data and code availability}

The complete run inventory (772 rows), per-run manifests, transcripts, evaluation outputs, annotation packs, the frozen codebook, adjudication records and all analysis scripts are available in the public repository mirror; the tables and figures in this article regenerate from the inventory by the released scripts (\texttt{p2\_reproduce\_tables.py}, \texttt{p2\_mechanism\_features.py}, \texttt{p2\_scan\_stats.py}, \texttt{p2c\_m0m1.py}). Total compute for all experiments reported here was approximately 767 agent runs at a cumulative API cost of roughly USD 105; local open-weights serving ran on a single RTX 4090. Internal planning documents, annotator-identity records and API credentials are excluded from the public mirror by design; the versioned repository history preserves the ordering of annotation freezing before unblinding as tamper-evident evidence of the pre-registration discipline.
